%%%%%%%%%%%%%%%%%%%%%%%%%%%%%%%%%%%%%%%%%%%%%%%%%%%%%%%%%%%%%%%%%%%%%%%%%%
%   FEHIME CEREN AY - SØKNADSBREV
%   Stilling: Senior analytiker, brukerinnsikt og måling
%   Nav, Arbeids- og velferdsdirektoratet (FINN-kode 470275492)
%   Søknadsfrist: 20. august 2026
%%%%%%%%%%%%%%%%%%%%%%%%%%%%%%%%%%%%%%%%%%%%%%%%%%%%%%%%%%%%%%%%%%%%%%%%%%

\documentclass[10pt,a4paper,roman]{moderncv}

\moderncvstyle{casual}
\moderncvcolor{blue}

\usepackage[utf8]{inputenc}
\usepackage[norsk]{babel}
\usepackage[scale=0.78,top=1.6cm,bottom=1.4cm]{geometry}
\setlength{\parskip}{0.9em}

%----------------------------------------------------------------------------------------
%   AVSENDER
%----------------------------------------------------------------------------------------

\firstname{Fehime Ceren}
\familyname{Ay}
\address{Bestumveien 37C, 0281}{Oslo}
\phone{+47 92 25 66 45}
\email{fcerenay@gmail.com}
\extrainfo{\href{https://cerenay.github.io}{\color{blue}cerenay.github.io}}

%----------------------------------------------------------------------------------------
%   MOTTAKER
%----------------------------------------------------------------------------------------

\recipient{Nav, Arbeids- og velferdsdirektoratet}{Avdeling brukeropplevelse, seksjon brukerinnsikt \\ Fyrstikkalléen 1 \\ 0661 Oslo}
\date{20. august 2026}
\opening{Søknad på stilling som senior analytiker, brukerinnsikt og måling}
\closing{Med vennlig hilsen,}
\enclosure[Vedlegg]{CV}

\begin{document}
\makelettertitle

Jeg søker med dette på stillingen som senior analytiker innen brukerinnsikt og
måling. Jeg har doktorgrad i økonomi og fem års erfaring med å bruke data til å
forstå hvordan tjenester faktisk brukes og oppleves, først i Telenor Research og
siden 2023 i Skatteetaten. Der arbeider jeg allerede med brukeropplevelse: jeg
kombinerer registerdata med atferdsdata fra Matomo og Skyra, i tverrfaglige team
sammen med utviklere og designere.

Det dere beskriver er den jobben jeg allerede gjør, i en etat med mange av de
samme utfordringene som Nav. Tidligere i år publiserte jeg og tre kolleger en analyse
av hvorfor noen ikke rapporterer utleieinntekter. Vi sendte en undersøkelse til
drøyt 18 000 personer, koblet svarene mot faktisk atferd i skattemeldingen og
mot tredjepartsdata, og brukte regresjon og segmentering til å skille fire
grupper med helt ulike årsaker. Det sterkeste enkeltfunnet var at digital
sårbarhet hadde klart størst negativ sammenheng med å rapportere riktig. Vi fant
også at en betydelig andel av dem som sa de kjente reglene, faktisk hadde
misforstått dem, og at moraliserende formuleringer gjorde folk mindre villige
til å lese informasjonen de trengte.

Jeg trekker fram nettopp denne studien fordi den handler om det Nav skriver i
annonsen: at alle skal kunne finne fram til, forstå og ta i bruk rettighetene
sine uavhengig av digitale ferdigheter og bakgrunn. Min erfaring er at et
etterlevelsesgap eller et bruksgap sjelden har én forklaring, og at nytten av
analysen ligger i å skille gruppene fra hverandre godt nok til at tiltakene kan
treffe ulikt.

Til daglig designer jeg spørreundersøkelser, eksperimenter og A/B-tester, og
måler effekten av endringer i tjenestene. Jeg er vant til å oversette det
designere og utviklere lurer på til noe som faktisk lar seg måle, og til å si fra
når datagrunnlaget ikke bærer konklusjonen de håper på. Analysene gjør jeg i SQL,
Python og R på Databricks. Jeg leverer også
styringsinformasjon og målekort til ledelsen, og bygger dashbord i Power BI og R
Shiny slik at fagmiljøene kan følge egne tall selv.

Fra Telenor har jeg erfaring med å etablere måling av brukertilfredshet for
digitale kanaler. Vi utviklet en metode for å predikere tilfredshet direkte fra
samtalelogger i kundeservice, slik at kvaliteten kunne følges løpende i stedet
for gjennom sporadiske undersøkelser. Arbeidet er publisert i Quality and User
Experience.

Doktorgraden min er i økonomi, med vekt på atferd og eksperimentelle metoder. Jeg er vant til å arbeide
hypotesedrevet og etterprøvbart, og til å være ærlig om hva dataene ikke kan
svare på. Det jeg motiveres av, er at analysene faktisk brukes. Å jobbe med
tjenester som betyr noe for folk i sårbare situasjoner er en vesentlig grunn til
at jeg søker.

Jeg stiller gjerne til samtale.

\makeletterclosing

\end{document}
